\documentclass[10pt,a4paper]{amsart}
\usepackage[margin=19mm]{geometry}
\usepackage{amsmath,amssymb,mathtools}
\usepackage[scr=rsfs,cal=euler]{mathalfa}
\usepackage{xfrac}
\usepackage[unicode,hidelinks,bookmarks=false]{hyperref}
\newcommand{\ie}{i.e.\ }
\newcommand{\cf}{cf.\ }
\newcommand{\resp}{respectively\ }
\newcommand{\Subsection}{\subsection}
\providecommand{\nobreakdash}{\nobreak\hbox{-}}
\setlength{\emergencystretch}{2em}
\multlinegap0pt
\usepackage{bm}
\usepackage{bbm}
\usepackage{dsfont}
\usepackage{xspace}
\usepackage{cancel}
\usepackage{mathtools}
\usepackage{amsthm}
\makeatletter
\@ifundefined{lemma}{\newtheorem{lemma}{Lemma}}{}
\makeatother
\makeatletter
\def\bee{\begin{equation}}
\def\ene{\end{equation}}
\def\mr#1{(\ref{#1})}
\expandafter\ifx\csname proof\endcsname\relax
\newenvironment{proof}{\begin{list}{$\!\!${\bf Proof.}%
  \rule{1pt}{0pt}}{\setlength{\leftmargin}{0pt}%
  \setlength{\itemindent}{30pt}%
  \setlength{\listparindent}{15pt}}\item}{\rule{0.3em}{0mm}%
  \hfill\framebox[1.2ex]{\rule{0.3em}{0mm}}\end{list}}
\fi
\expandafter\ifx\csname prooft\endcsname\relax
\newenvironment{prooft}[1]{\begin{list}{$\!\!${\bf Proof of Theorem~\ref{#1}.}%
  \rule{1pt}{0pt}}{\setlength{\leftmargin}{0pt}%
  \setlength{\itemindent}{30pt}%
  \setlength{\listparindent}{15pt}}\item}{\rule{0.3em}{0mm}%
  \hfill\framebox[1.2ex]{\rule{0.3em}{0mm}}\end{list}}
\fi
\makeatother
\title[Cubic spline functions revisited]{Cubic spline functions revisited}
\author{Florian Jarre}
\date{Source: arXiv:2507.05083v1 (CC BY 4.0)}
\begin{document}
\maketitle

\noindent\emph{Source and licence.} Cubic spline functions revisited. Version arXiv:2507.05083v1 (CC BY 4.0), distributed under the Creative Commons Attribution 4.0 International licence, as recorded in the arXiv licence metadata for this exact version and shown on the abstract page https://arxiv.org/abs/2507.05083v1. Full rights notice: licenses/arxiv-2507.05083-CC-BY-4.0.txt.\par\smallskip

\noindent\textit{Editorial scope.} Counted unit: the Lemma whose source label is \texttt{3128} in the article (source0.tex lines 194--210, source label \texttt{3128}), delivered together with the article's own complete original proof of it (source0.tex lines 213--243, one \texttt{proof} environment of 31 lines). No mathematics is written, repaired or re-derived: statement and proof are the article's own text line for line. Required context retained uncounted: lemma1 (lines 135--142). Every definition, display and label of those lines is retained unchanged. Omitted from this package and not redistributed: 1--134, 143--193, 211--212, 244--1340. Nothing inside a retained excerpt is omitted or altered: every retained body is delivered exactly as the article's lines are written, so no omission receipt of any kind accompanies this package. Citations: the retained excerpts carry no citation command at all, so no bibliography entry of the article is transcribed and no reference is redistributed. Reference adaptation: none was needed and none was made; every reference inside the delivered unit resolves inside it, so no unresolved reference or citation is produced. Printed slips kept literal: no printed word, symbol, hypothesis or number of the counted statement or of its proof was altered. Layout and class: the article's own class is \texttt{article} with packages including inputenc, graphicx, amsmath, amssymb, color, url; the wrapper is the lane's pinned \texttt{amsart} editorial wrapper at 10pt on A4, which restates the theorem-like environments the retained text needs and adds the article's own macro definitions that the retained text needs (\texttt{\textbackslash bee}, \texttt{\textbackslash ene}, \texttt{\textbackslash mr}), with extra wrapper packages (bm, bbm, dsfont, xspace, cancel, mathtools). The delivered numbering is the unit's own. Assets: no retained span contains a figure environment, a graphics-inclusion command, a PDF-page inclusion, a table or a TikZ picture; no image file is copied into this package, the package declares no asset dependency and no image file is redistributed with it. Licence: version arXiv:2507.05083v1 of this article is distributed under the Creative Commons Attribution 4.0 International licence, as recorded in the arXiv licence metadata for this exact version and shown on the abstract page https://arxiv.org/abs/2507.05083v1. A full rights notice is delivered at licenses/arxiv-2507.05083-CC-BY-4.0.txt. Characters: every retained excerpt is pure ASCII, so the delivered engine, XeLaTeX with its default Latin Modern text font, reports no missing character.
\begin{lemma}\label{lemma1}
  The cubic spline function with either the first derivative
  or the second derivative specified at both end points 
  is the unique function that
  minimizes the integral of the square
  of the second derivative among all interpolating $C^2$-functions
  with the same end conditions.
\end{lemma}

\begin{lemma}\label{3128}
  Let $f$ be four times continuously differentiable on $[x_0,x_n]$.
  Based on the data $f(x_i)$ for $0\le i\le n$,
  %as well as $f'(x_0)$ and $f'(x_n)$,
  an approximation of $f(x)$ for $x\in[x_1,x_{n-1}]$ is possible
  by an interpolating cubic polynomial $p$ such that 
  $$
  |p(x)-f(x)|\le \tfrac{3}{128}\|f^{(4)}\|_\infty h^4
  $$
  with $h:=\max_{1\le i\le n}(x_i-x_{i-1})$.
  
  For $x\in[x_0,x_1]$ or $x\in[x_{n-1},x_n]$ 
  $$
  |p(x)-f(x)|\le \tfrac{1}{24}\|f^{(4)}\|_\infty h^4
  $$
  Both estimates are  best possible.
\end{lemma}

\begin{proof}
  When $x$ coincides with $x_i$ for some $i$ there is nothing to show.
  Let $x\in(x_1,x_{n-1})$. Choose $i$ such that $x\in(x_i,x_{i+1})$ and
  let $p$ be the cubic polynomial interpolating $f$ at
  $x_{i-1},x_i,x_{i+1},x_{i+2}$.
  The standard error bound
  for polynomial interpolation states that
  \bee\label{poly1}
  f(x)-p(x)=\tfrac{f^{(4)}(\xi)}{24} \omega(x) \quad\hbox{with}
  \quad \omega(x) = \prod_{j=i-1}^{i+2}(x-x_j)
  \ene
  and with $\xi\in (x_{i-1},x_{i+2})$.
  Observe that $\|\omega\|_\infty$ increases when, for example,
  $x_{i-1}$ is reduced while $x_i, x_{i+1}$, and $x_{i+2}$ remain
  unchanged. Thus, when maximizing $\|\omega\|_\infty$ one can
  assume without loss of generality that all mesh points have
  maximum distance $x_{k+1}-x_k = h$ for all $k$.
  Straightforward calculations then show that
  $|\omega(x)|\le \tfrac 9{16}h^4$
  for $x\in [x_i,x_{i+1}]$.
  The first statement of the 
  lemma follows when inserting this into \mr{poly1}.

  Similarly, for $x\in [x_0,x_{1}]$, the term $\omega(x)$ is given by
  $\omega(x)= \prod_{j=0}^{4}(x-x_j)$ and straightforward calculations show that
  $|\omega(x)|\le h^4$.
  Likewise also for $x\in[x_{n-1},x_n]$.

  When interpolating the function $f$ with $f(x)\equiv x^4$, the fourth derivative is constant
  so that \mr{poly1} implies that the bounds given in Lemma \ref{lemma1} are best possible.
\end{proof}

\end{document}
